%% toc: P2 19 19.1,19.2,19.3
\section{Herbrand structures and ground satisfaction}
\label{sec:p02-w02-p02-028-herbrand-structures}

This section and the next prove Herbrand's theorem in its validity form and
connect it with the clause form and with resolution.  The definitions of the
Herbrand universe, Herbrand base, Herbrand structure and finite Herbrand
expansion are those of
Definitions~\ref{def:p02-herbrand-universe-base},
\ref{def:p02-herbrand-structure} and~\ref{def:p02-herbrand-expansion}; the
statement proved first is Theorem~\ref{thm:p02-herbrand-validity}.

\emph{Purpose.}  Herbrand's theorem reduces first-order validity of a sentence
to propositional validity of one finite disjunction of ground instances of a
quantifier-free formula.  It is the semantic basis of ground and lifted
resolution and of Statman's analysis of direct computations in
Section~\ref{sec:p02-w02-p02-028-statman}.
\emph{Inputs.}  An equality-free first-order signature \(\Lang\) and an
\(\Lang\)-sentence; for the clause form, a finite equality-free clause set.
\emph{Conclusions.}  Theorems~\ref{thm:p02-w02-p02-028-existential-validity}
and~\ref{thm:p02-w02-p02-028-general-validity} (validity form),
Theorem~\ref{thm:p02-w02-p02-028-clause-form} (clause form),
Theorem~\ref{thm:p02-w02-p02-028-validity-to-clauses} (the passage between
them), and Theorem~\ref{thm:p02-w02-p02-028-resolution} (the passage to
resolution).  All proofs are in the annex.

\begin{hbconvention}{def:p02-w02-p02-028-standing}{Standing assumptions}
Throughout Sections~\ref{sec:p02-w02-p02-028-herbrand-structures}
and~\ref{sec:p02-w02-p02-028-validity-forms}:
\begin{enumerate}
  \item \emph{Equality.}  Formulas and clauses are equality-free: no atom has
        the form \(r=t\).  Logical equality is not treated as an ordinary
        binary relation symbol, so the results below say nothing about
        sentences containing it.
  \item \emph{Ground terms.}  \(\Lang\) contains a constant symbol, hence a
        ground term.  When a given signature has none, one fresh constant is
        added first; the paragraph before
        Definition~\ref{def:p02-herbrand-universe-base} shows that this
        changes neither validity nor satisfiability of sentences of the old
        signature.
  \item \emph{Structures} have nonempty domains, as everywhere in this part.
  \item \emph{Choice.}  The Axiom of Choice is used only where it is named:
        in the expansion direction of Skolemization
        (Theorem~\ref{thm:skolem-equisatisfiable}).
\end{enumerate}
\end{hbconvention}

\begin{hbdefinition}{def:p02-w02-p02-028-atom-valuation}{Atom valuations and propositional truth of ground formulas}
An \emph{atom valuation} is a map \(v\) from ground atoms of \(\Lang\) to
\(\{0,1\}\).  For a ground quantifier-free formula \(\psi\), its propositional
value \(\llbracket\psi\rrbracket_v\) is defined by \(\llbracket A\rrbracket_v=v(A)\)
for ground atoms, \(\llbracket\bottom\rrbracket_v=0\), and the truth-table
clauses for \(\neg,\land,\lor,\to\).  The formula is \emph{propositionally
valid} when its value is \(1\) under every atom valuation, and a set of such
formulas is \emph{propositionally satisfiable} when one atom valuation gives
every member the value \(1\).  This is the reading fixed in
Definition~\ref{def:p02-herbrand-expansion}: each distinct ground atom is one
propositional variable.  For a structure \(\mathcal N\), its \emph{induced atom
valuation} is \(v_{\mathcal N}(A)=1\) iff \(\mathcal N\entails A\).  For a set
\(B\) of ground atoms, \(v_B\) is the characteristic function of \(B\).
\end{hbdefinition}

\begin{hblemma}{thm:p02-w02-p02-028-ground-evaluation}{proved=annex:p02-w02-p02-028-ground-evaluation}{Ground evaluation}
Let \(\mathcal N\) be an \(\Lang\)-structure.
\begin{enumerate}
  \item For every ground quantifier-free \(\psi\),
        \(\mathcal N\entails\psi\) iff \(\llbracket\psi\rrbracket_{v_{\mathcal N}}=1\).
  \item For every quantifier-free \(\varphi(x_1,\ldots,x_n)\), ground terms
        \(t_1,\ldots,t_n\) and assignment \(s\),
        \(\mathcal N,s[x_1\mapsto t_1^{\mathcal N},\ldots,x_n\mapsto t_n^{\mathcal N}]\entails\varphi\)
        iff \(\mathcal N\entails\varphi(t_1,\ldots,t_n)\), where
        \(\varphi(t_1,\ldots,t_n)\) is the simultaneous replacement of each
        \(x_i\) by \(t_i\).
  \item If \(\mathcal N=\mathcal M_B\) is a Herbrand structure, then
        \(t^{\mathcal M_B}=t\) for every ground term \(t\), and
        \(v_{\mathcal M_B}=v_B\).
\end{enumerate}
\end{hblemma}

\begin{hblemma}{thm:p02-w02-p02-028-herbrand-model}{proved=annex:p02-w02-p02-028-herbrand-model}{Herbrand model lemma}
Let \(\Lang\) be countable, and let \(\Phi\) be a set of quantifier-free
\(\Lang\)-formulas.  Let \(\Gamma\) be the set of all ground instances
\(\varphi(t_1,\ldots,t_n)\) of members \(\varphi(x_1,\ldots,x_n)\) of
\(\Phi\), where \(x_1,\ldots,x_n\) are the variables occurring in
\(\varphi\) and the \(t_i\) range over the Herbrand universe of \(\Lang\).
If every finite subset of \(\Gamma\) is propositionally satisfiable, then
there is a set \(B\) of ground atoms such that the Herbrand structure
\(\mathcal M_B\) satisfies the universal closure of every member of
\(\Phi\).
\end{hblemma}

Countability is used only to apply propositional compactness
(Theorem~\ref{thm:prop-compactness}), whose propositional variables form the
countable set \(\{p_n:n\in\mathbb N\}\).  Every application below first passes
to the finite signature of the symbols actually occurring in the input, so
no restriction on the size of \(\Lang\) itself is imposed by the theorems.

\section{Validity forms of Herbrand's theorem}
\label{sec:p02-w02-p02-028-validity-forms}

\subsection{The existential case}
\label{sec:p02-w02-p02-028-existential}

\begin{hbtheorem}{thm:p02-w02-p02-028-existential-validity}{proved=annex:p02-w02-p02-028-existential-validity}{Herbrand's theorem for existential sentences}
Let \(\sigma=\exists x_1\cdots\exists x_n\,\varphi(x_1,\ldots,x_n)\) be an
equality-free \(\Lang\)-sentence with quantifier-free \(\varphi\).  Then
\(\sigma\) is valid iff some finite Herbrand expansion
\(\varphi(\bar t_1)\lor\cdots\lor\varphi(\bar t_m)\) of \(\sigma\), with
\(m\geq1\), is propositionally valid.  Moreover, the ground terms in the
tuples \(\bar t_i\) can be chosen to contain only the function and constant
symbols occurring in \(\sigma\), together with one constant of \(\Lang\) when
\(\sigma\) contains none.
\end{hbtheorem}

This is the statement of Theorem~\ref{thm:p02-herbrand-validity}, with the
additional information on the symbols in the instance terms; the annex proof
proves both directions.  The right-to-left direction uses no compactness; the
left-to-right direction uses Lemma~\ref{thm:p02-w02-p02-028-herbrand-model}.
A single instance need not suffice (Example~\ref{def:p02-w02-p02-028-drinker}).

\subsection{Herbrand normal form and the general validity form}
\label{sec:p02-w02-p02-028-herbrand-normal-form}

\begin{hbdefinition}{def:p02-w02-p02-028-herbrand-form}{Herbrand form}
Let \(\sigma=Q_1x_1\cdots Q_nx_n\,\theta\) be a rectified prenex
\(\Lang\)-sentence with quantifier-free \(\theta\).  Traverse the prefix from
left to right.  At each universal occurrence \(\forall x_i\), list the
variables \(u_1,\ldots,u_m\) bound by the existential occurrences to its
left, introduce a new \(m\)-ary function symbol \(h_i\) (a constant when
\(m=0\)), replace the occurrences of \(x_i\) in the matrix by
\(h_i(u_1,\ldots,u_m)\), and delete \(\forall x_i\).  Retain every
existential.  Distinct universal occurrences receive distinct new symbols.
The result is an existential sentence
\[
  \operatorname{Hf}(\sigma)=\exists u_1\cdots\exists u_p\,\theta^{\mathrm H}
\]
in the signature \(\Lang^{\mathrm H}\) obtained by adding the new symbols,
where \(u_1,\ldots,u_p\) are the existential variables in prefix order.  It is
the dual of Skolemization (Definition of \(\operatorname{Sk}\) before
Theorem~\ref{thm:skolem-equisatisfiable}): the roles of the two quantifiers
are exchanged.
\end{hbdefinition}

\begin{hblemma}{thm:p02-w02-p02-028-herbrand-skolem-duality}{proved=annex:p02-w02-p02-028-duality}{Herbrand form as negated Skolem form}
For \(\sigma\) as in Definition~\ref{def:p02-w02-p02-028-herbrand-form}, let
\(\sigma^{\neg}=\overline{Q}_1x_1\cdots\overline{Q}_nx_n\,\neg\theta\), where
\(\overline{\forall}=\exists\) and \(\overline{\exists}=\forall\).  Then
\(\sigma^{\neg}\) is a rectified prenex sentence logically equivalent to
\(\neg\sigma\), and, when the Skolem symbol of \(\sigma^{\neg}\) at position
\(i\) is taken to be \(h_i\),
\[
 \operatorname{Sk}(\sigma^{\neg})=\forall u_1\cdots\forall u_p\,\neg\theta^{\mathrm H}.
\]
\end{hblemma}

\begin{hbtheorem}{thm:p02-w02-p02-028-herbrand-form-validity}{proved=annex:p02-w02-p02-028-hf-validity}{Validity is preserved by the Herbrand form}
Let \(\sigma\) be a rectified prenex equality-free \(\Lang\)-sentence.  If
\(\sigma\) is valid, then \(\operatorname{Hf}(\sigma)\) is valid in
\(\Lang^{\mathrm H}\); this direction uses no choice principle.  Assuming the
Axiom of Choice, if \(\operatorname{Hf}(\sigma)\) is valid, then \(\sigma\) is
valid.
\end{hbtheorem}

\begin{hbtheorem}{thm:p02-w02-p02-028-general-validity}{proved=annex:p02-w02-p02-028-general-validity}{Herbrand's theorem, validity form for arbitrary sentences}
Assume the Axiom of Choice.  Let \(\sigma\) be any equality-free
\(\Lang\)-sentence, let \(\sigma'\) be a rectified prenex sentence logically
equivalent to it (Theorem~\ref{thm:prenex-normal-form}), and let
\(\operatorname{Hf}(\sigma')=\exists\bar u\,\theta^{\mathrm H}\).  Then
\(\sigma\) is valid iff there are finitely many tuples
\(\bar t_1,\ldots,\bar t_m\) (\(m\geq1\)) of ground
\(\Lang^{\mathrm H}\)-terms such that
\[
 \theta^{\mathrm H}(\bar t_1)\lor\cdots\lor\theta^{\mathrm H}(\bar t_m)
\]
is propositionally valid.  The implication from validity of \(\sigma\) to the
existence of such a disjunction uses no choice principle.
\end{hbtheorem}

The ground instances in the general form are instances over the Herbrand
universe of the expanded signature \(\Lang^{\mathrm H}\), not of \(\Lang\): the
new symbols \(h_i\) name the values at which the universal quantifiers of
\(\sigma'\) must be checked.  When \(\sigma\) is already existential, no new
symbol is introduced and Theorem~\ref{thm:p02-w02-p02-028-general-validity}
reduces to Theorem~\ref{thm:p02-w02-p02-028-existential-validity}.

\begin{hbexample}{def:p02-w02-p02-028-drinker}{One instance is not enough}
Let \(P\) be unary and \(\sigma=\exists x\,\forall y\,(P(x)\to P(y))\).  Its
Herbrand form is \(\exists x\,(P(x)\to P(h(x)))\) with a new unary \(h\); add
a constant \(c\) because the signature has none.  No single instance
\(P(t)\to P(h(t))\) is propositionally valid: give \(P(t)\) the value \(1\)
and \(P(h(t))\) the value \(0\).  The two-member expansion
\[
 \bigl(P(c)\to P(h(c))\bigr)\lor\bigl(P(h(c))\to P(h(h(c)))\bigr)
\]
is propositionally valid: if the first disjunct has value \(0\), then
\(P(h(c))\) has value \(0\), so the second disjunct has value \(1\).  By
Theorem~\ref{thm:p02-w02-p02-028-general-validity}, \(\sigma\) is valid.
\end{hbexample}

\begin{hbwarning}{def:p02-w02-p02-028-equality-limit}{Equality and recursive bounds are outside these theorems}
With logical equality the propositional reading of ground atoms is wrong: the
formula \(\exists x\,(x=c)\) is valid, but the ground instance \(c=c\) is not
true under every atom valuation when \(c=c\) is treated as an independent
propositional variable.  An equality version needs equality axioms among the
instances or an equality-aware calculus, as noted after
Theorem~\ref{thm:fol-resolution}.  The theorems also give no recursive bound
on the number \(m\) of instances in terms of \(\sigma\) alone: such a bound
would make validity decidable, contradicting
Theorem~\ref{thm:fol-validity-undecidable}
(Corollary~\ref{thm:p02-w02-p02-028-no-recursive-bound}).
\end{hbwarning}

\subsection{From the validity form to clauses and resolution}
\label{sec:p02-w02-p02-028-clauses}

\begin{hbtheorem}{thm:p02-w02-p02-028-clause-form}{proved=annex:p02-w02-p02-028-clause-form}{Herbrand's theorem for clauses}
Let \(S\) be a finite equality-free set of first-order clauses in \(\Lang\).
Then \(S\) is unsatisfiable iff some finite set of ground instances of
members of \(S\), read as propositional clauses whose distinct ground atoms
are propositional variables, is propositionally unsatisfiable.  The ground
terms can be chosen from the symbols occurring in \(S\), together with one
constant of \(\Lang\) when \(S\) contains none.
\end{hbtheorem}

This is the statement of Theorem~\ref{thm:herbrand-clauses}, proved here in
both directions.

\begin{hbtheorem}{thm:p02-w02-p02-028-validity-to-clauses}{proved=annex:p02-w02-p02-028-validity-to-clauses}{Validity form and clause form}
Let \(\sigma\) be an equality-free \(\Lang\)-sentence, \(\sigma'\) and
\(\operatorname{Hf}(\sigma')=\exists\bar u\,\theta^{\mathrm H}\) as in
Theorem~\ref{thm:p02-w02-p02-028-general-validity}, and let
\[
 S_\sigma=\operatorname{Cl}\bigl(\operatorname{nnf}(\neg\theta^{\mathrm H})\bigr)
\]
be the clause family of Method~\ref{def:p02-matrix-clause-construction}, read
universally.  Then:
\begin{enumerate}
  \item for ground tuples \(\bar t_1,\ldots,\bar t_m\) of
        \(\Lang^{\mathrm H}\)-terms, the disjunction
        \(\theta^{\mathrm H}(\bar t_1)\lor\cdots\lor\theta^{\mathrm H}(\bar t_m)\) is
        propositionally valid iff the set of ground clauses
        \(\{C(\bar t_j):C\in S_\sigma,\ 1\leq j\leq m\}\) is propositionally
        unsatisfiable, where \(C(\bar t_j)\) replaces each \(u_i\) by the
        \(i\)-th member of \(\bar t_j\);
  \item every finite set of ground instances of members of \(S_\sigma\) is
        contained in a set of the form in item~1; and
  \item assuming the Axiom of Choice, \(\sigma\) is valid iff \(S_\sigma\) is
        unsatisfiable; the implication from validity of \(\sigma\) to
        unsatisfiability of \(S_\sigma\) uses no choice principle.
\end{enumerate}
\end{hbtheorem}

Thus the finite Herbrand expansions of the validity form and the finite
unsatisfiable sets of ground instances of the clause form are the same
objects read dually: a valid disjunction of instances of
\(\theta^{\mathrm H}\) is a refutable conjunction of instances of its
negation, and negation followed by Skolemization is exactly passage to the
Herbrand form (Lemma~\ref{thm:p02-w02-p02-028-herbrand-skolem-duality}).

The lifting step below needs most general unifiers.  The first-order
unification procedure of \cref{sec:fol-cardinality-normalforms-04} is used as
defined there: it starts from the equations pairing corresponding arguments
of the two atoms, and repeatedly deletes \(t\doteq t\), decomposes
\(f(r_1,\ldots,r_n)\doteq f(t_1,\ldots,t_n)\), fails on a clash of function
symbols or arities, orients \(t\doteq x\) to \(x\doteq t\), and, for
\(x\doteq t\) with \(x\neq t\), either fails by the occurs check or eliminates
\(x\) by \(\{x\mapsto t\}\), composing that substitution onto an accumulator.

\begin{hblemma}{thm:p02-w02-p02-028-unification}{proved=annex:p02-w02-p02-028-unification}{Termination and correctness of unification}
Let \(A\) and \(B\) be atoms, and run the first-order unification procedure
on them, selecting the equation to be treated at each step in any way.
\begin{enumerate}
  \item The run terminates.
  \item It ends in failure iff \(A\) and \(B\) have no unifier.
  \item If it ends with the equation set empty, the accumulated substitution
        \(\theta\) is a most general unifier of \(A\) and \(B\).
\end{enumerate}
In particular, any two unifiable atoms have a most general unifier.
\end{hblemma}

This is the theorem on first-order unification of
\cref{sec:fol-cardinality-normalforms-04}, with a complete proof: the annex
gives the termination measure and checks the invariant for each
transformation.

\begin{hblemma}{thm:p02-w02-p02-028-lifting}{proved=annex:p02-w02-p02-028-lifting}{Lifting of ground resolution}
Use the first-order Resolution and Factoring rules and renaming apart of
\cref{sec:fol-cardinality-normalforms-04}.  Let \(S\) be a set of first-order
clauses and let \(K_1,\ldots,K_r\) be a propositional resolution derivation
(\cref{sec:prop-metatheory-04}) in which every \(K_j\) is either a ground
instance of a member of \(S\) or a resolvent of two earlier members.  Then for
every \(j\) there are a first-order resolution derivation from \(S\) with
factoring that contains a clause \(K_j'\), and a substitution \(\gamma_j\)
sending the variables of \(K_j'\) to ground terms, with
\(K_j'\gamma_j=K_j\).  In particular, if \(K_r=\varnothing\), then the empty
clause has a first-order resolution derivation from \(S\).
\end{hblemma}

The lemma uses, for each Factoring and Resolution step, the existence of a
most general unifier of two unifiable atoms, which is
Lemma~\ref{thm:p02-w02-p02-028-unification}.

\begin{hbtheorem}{thm:p02-w02-p02-028-resolution}{proved=annex:p02-w02-p02-028-resolution}{From the clause form to resolution}
Let \(S\) be a finite equality-free set of first-order clauses.  The
following are equivalent:
\begin{enumerate}
  \item \(S\) is unsatisfiable;
  \item some finite set of ground instances of members of \(S\) has a
        propositional resolution refutation;
  \item first-order Resolution with Factoring derives the empty clause from
        \(S\).
\end{enumerate}
Consequently, for an equality-free sentence \(\sigma\), assuming the Axiom of
Choice, \(\sigma\) is valid iff Resolution with Factoring derives the empty
clause from \(S_\sigma\) of
Theorem~\ref{thm:p02-w02-p02-028-validity-to-clauses}.
\end{hbtheorem}

This proves the forward implication of Theorem~\ref{thm:fol-resolution}
along the route its sketch describes, with the clause theorem, the
unification theorem and the lifting step now proved.

\begin{hbcorollary}{thm:p02-w02-p02-028-no-recursive-bound}{proved=annex:p02-w02-p02-028-no-recursive-bound}{Search for instances and the absence of a recursive bound}
Fix an effective coding of first-order formulas over a countable signature.
\begin{enumerate}
  \item There is a partial recursive procedure that, applied to an
        equality-free existential sentence \(\exists\bar x\,\varphi\), halts
        exactly when the sentence is valid, and then outputs tuples
        \(\bar t_1,\ldots,\bar t_m\) with
        \(\varphi(\bar t_1)\lor\cdots\lor\varphi(\bar t_m)\) propositionally
        valid.
  \item If equality-free validity is undecidable, then there is no recursive
        function that assigns to every valid equality-free existential
        sentence a finite list of tuples whose expansion is propositionally
        valid, and is defined on all equality-free existential sentences.
\end{enumerate}
\end{hbcorollary}

\section{Statman: Herbrand's theorem and Gentzen's notion of a direct proof}
\label{sec:p02-w02-p02-028-statman}

This section presents R.~Statman's chapter of the same title
\cite[pp.~897--912]{Statman1977}.  Numbering in square brackets such as
[2.10] is the chapter's own; statements and notation follow the chapter.
Where the chapter gives a proof, the annex writes it out in full; where the
chapter only indicates a route or leaves details to the reader, the annex
follows that route and supplies the missing steps, and the reader text says
where a step or a hypothesis is the handbook's rather than the chapter's.
The chapter calls the Horn statement of [1.3(ii)] a simple consequence of
Gentzen's cut-elimination theorem and gives no derivation.
Theorem~\ref{thm:p02-w02-p02-028-st-horn-cut} carries out that derivation,
from the cut-elimination theorem for the defined-negation sequent calculus,
and states the hypotheses the route uses; Proposition~\ref{thm:p02-w02-p02-028-st-horn}
proves the same statement a second way, through Herbrand's theorem.

\emph{Purpose.}  The chapter uses Herbrand's theorem as a point of departure
and asks finer questions about one restricted class of formulas: which
substitution instances are needed, how proof length depends on them, and in
what order rules should be applied.  It answers them for the equation
calculus \(K\) of Kreisel and Tait.
\emph{Inputs.}  A finite set \(S\) of equations between individual terms
built from individual variables, \(0\), \(s\) and function variables, and an
equation \(E\).
\emph{Conclusions.}  Completeness and decidability of \(K\); bounds on the
length of direct computations
(Propositions~\ref{thm:p02-w02-p02-028-st-prop-3-1}
and~\ref{thm:p02-w02-p02-028-st-prop-3-2}); a polynomial-time decision
procedure (Theorem~\ref{thm:p02-w02-p02-028-st-decision}); and the
permutation analysis of singular computations
(Propositions~\ref{thm:p02-w02-p02-028-st-prop-3-3-2}
to~\ref{thm:p02-w02-p02-028-st-prop-4-1-1}).

\subsection{The setting: Herbrand's theorem and restricted classes}
\label{sec:p02-w02-p02-028-st-setting}

\emph{[1.1] Herbrand's theorem as the starting point.}  The chapter states, as
a corollary of the completeness theorem, that if \(G(x_1,\ldots,x_n)\) is
quantifier-free and \(\exists x_1\cdots x_nG(x_1,\ldots,x_n)\) is valid, then
there are terms \(t^i_j\) (\(1\leq i\leq m\), \(1\leq j\leq n\)) such that
\(\bigvee_{1\leq i\leq m}G(t^i_1,\ldots,t^i_n)\) is valid
\cite[\S1.1, p.~898]{Statman1977}.  It refers for this result to Theorem~2.9
of Chapter~D.2 of the same handbook \cite[Theorem~2.9, p.~876]{Schwichtenberg1977}.
Here \(G\) is any quantifier-free formula of predicate logic, so it may
contain the equality sign, and the conclusion is validity, not propositional
validity.  Theorem~\ref{thm:p02-w02-p02-028-st-herbrand-eq} proves the
statement in that generality.  For equality-free \(G\) the left-to-right
direction of Theorem~\ref{thm:p02-w02-p02-028-existential-validity} gives the
stronger conclusion of propositional validity; a propositionally valid ground
formula is valid by Lemma~\ref{thm:p02-w02-p02-028-ground-evaluation}.

\begin{hbtheorem}{thm:p02-w02-p02-028-st-herbrand-eq}{proved=annex:p02-w02-p02-028-st-herbrand-eq}{Herbrand's theorem with equality [1.1]}
Let \(G(x_1,\ldots,x_n)\) be a quantifier-free formula of a first-order
signature \(\Lang\), possibly containing \(=\) and free variables other than
\(x_1,\ldots,x_n\).  If \(\exists x_1\cdots\exists x_n\,G\) is valid, then
there are \(m\geq1\) and \(\Lang\)-terms \(t^i_j\)
(\(1\leq i\leq m\), \(1\leq j\leq n\)) such that
\(\bigvee_{1\leq i\leq m}G(t^i_1,\ldots,t^i_n)\) is valid.  Conversely, if
such a disjunction is valid, so is \(\exists x_1\cdots\exists x_n\,G\).
\end{hbtheorem}

The proof reads \(=\) as a new binary relation symbol \(\approx\), adds the
equality axioms for the symbols of \(G\) as universal hypotheses, applies the
existential case of Herbrand's theorem to the resulting equality-free
sentence, and discards the instances of the equality axioms, which are true
when \(\approx\) is read as identity.

\emph{[1.2] Methods of finding the terms.}  The chapter observes that the terms
and a derivation of the disjunction can be obtained recursively from any
derivation of \(\exists\bar x\,G\); that a partial recursive method working
from the formula alone exists; and that no recursive method working from the
formula alone exists, because validity is undecidable
\cite[\S1.2, p.~898]{Statman1977}.  For the equality-free case these three
claims are proved in Corollary~\ref{thm:p02-w02-p02-028-no-recursive-bound}
and the annex note to it; for the general case they are
Corollary~\ref{thm:p02-w02-p02-028-st-methods}.

\begin{hbcorollary}{thm:p02-w02-p02-028-st-methods}{proved=annex:p02-w02-p02-028-st-methods}{Finding the terms [1.2]}
Fix an effective coding of first-order formulas and of derivations in the
first-order natural deduction calculus of Theorems~\ref{thm:fol-soundness}
and~\ref{thm:fol-completeness}, and assume the Axiom of Choice.
\begin{enumerate}
  \item Validity of quantifier-free formulas, with or without equality, is
        decidable.
  \item There is a partial recursive procedure that, applied to a formula
        \(\exists\bar x\,G\) with \(G\) quantifier-free, halts exactly when
        the formula is valid, and then outputs terms \(t^i_j\) as in
        Theorem~\ref{thm:p02-w02-p02-028-st-herbrand-eq} together with a
        natural deduction derivation of the disjunction.
  \item There is a total recursive procedure that does the same when applied
        to any natural deduction derivation of \(\exists\bar x\,G\) without
        open assumptions.
  \item No total recursive procedure assigns to every formula
        \(\exists\bar x\,G\) a list of terms that makes the disjunction valid
        whenever \(\exists\bar x\,G\) is valid.
\end{enumerate}
Item~4 is proved for a coding whose signature has infinitely many function
symbols of each arity, which the Herbrand form of
Definition~\ref{def:p02-w02-p02-028-herbrand-form} draws on; it uses the
undecidability of first-order validity
(Theorem~\ref{thm:fol-validity-undecidable}), which the handbook cites, for
that coding.
\end{hbcorollary}

\emph{[1.3] Three refinements.}  The theory of proofs, as the chapter
describes it, restricts attention in three ways.
\begin{enumerate}
  \item[(i)] \emph{Restricted formula classes.}  Fix prime formulas
        \(A_1,\ldots,A_p\) and implicational formulas \(I_1,\ldots,I_q\) with
        \(I_j=P^j_1\to(\cdots(P^j_{r(j)}\to C^j)\cdots)\) for prime
        \(P^j_i\) and \(C^j\).  The class \(L\) consists of the formulas
        \(G\) of the form
        \[
          \Bigl(\bigwedge_{1\leq i\leq p}A_i\land\bigwedge_{1\leq j\leq q}I_j
          \land\bigwedge_{1\leq k\leq m}B_k\Bigr)\to B_{m+1}
        \]
        with \(B_k\) prime and not containing \(x_1,\ldots,x_n\)
        \cite[\S1.3(i), p.~898]{Statman1977}.  The motivating instance is the
        language of semigroups, with formulas saying that two words are equal
        in a finitely generated, finitely presented cancellation semigroup.
  \item[(ii)] \emph{Rules adapted to \(L\).}  For relations
        \(R_1,\ldots,R_m\) of a presentation, if \(w_1=w_2\) holds then it
        is derivable from \(R_1,\ldots,R_m\), the associativity axioms
        \(((ab)c)=(a(bc))\) and the identity axioms \(a=a\) by substituting
        equals for equals and the two cancellation rules
        \[
          \frac{(ab)=(ac)}{b=c}
          \qquad\text{and}\qquad
          \frac{(ba)=(ca)}{b=c}.
        \]
        In general, if \(\exists\bar x\,G\) is valid for \(G\in L\), then
        \(B_{m+1}\) is derivable from \(B_1,\ldots,B_m\) and the axioms
        \(\theta A_i\), for substitutions \(\theta\) for \(x_1,\ldots,x_n\),
        by the rules
        \[
          \frac{\theta P^j_1\quad\cdots\quad\theta P^j_{r(j)}}{\theta C^j};
        \]
        the chapter calls this a simple consequence of Gentzen's
        cut-elimination theorem \cite[\S1.3(ii), pp.~898--899]{Statman1977}.
        Both claims are proved under explicit hypotheses in
        Propositions~\ref{thm:p02-w02-p02-028-st-semigroup}
        and~\ref{thm:p02-w02-p02-028-st-horn}; the second is also derived
        from cut elimination, as the chapter indicates, in
        Theorem~\ref{thm:p02-w02-p02-028-st-horn-cut}.
  \item[(iii)] \emph{Length as well as term complexity.}  The chapter raises
        three questions \cite[\S1.3(iii), p.~899]{Statman1977}:
        (a) which substitutions \(\theta\) are needed to derive \(B_{m+1}\)
        from \(B_1,\ldots,B_m\);
        (b) how the length of such a derivation is related to the complexity
        of the \(\theta\) used; and
        (c) what is an efficient order in which to apply the rules.
        It stresses that these questions remain substantive even when
        validity in \(L\) is decidable, since they concern the efficiency of
        decision procedures.
\end{enumerate}

\begin{hbdefinition}{def:p02-w02-p02-028-st-direct}{Direct proof [1.4]}
In the chapter's informal sense, a \emph{direct} proof is one containing no
term more complicated than all the terms of the theorem proved
\cite[\S1.4, p.~899]{Statman1977}.  The chapter attributes to the theory of
proofs the aim of making such differences among proofs objects of study, and
quotes the phrase ``aesthetic criteria of elegance or convenience'' for the
older standard.  Its exact replacement for ``more complicated'' is a
measure (Definition~\ref{def:p02-w02-p02-028-st-measure}), and the exact notion
is \(M\)-directness (Definition~\ref{def:p02-w02-p02-028-st-m-direct}).
\end{hbdefinition}

If direct proofs can be formalized in a system of rules for \(L\), this
answers question (a); the chapter notes that such proofs can then be the
outputs of a complete proof-search procedure resembling semantic tableaux.

\emph{[1.4] The Kreisel--Tait calculus and the axioms \(J\).}  In 1961
Kreisel and Tait gave an equation calculus \(K\) meant to formalize
computation from equations that define the values of a number-theoretic
function \cite{KreiselTait1961}.  The chapter relates \(K\) to a set \(J\) of
axioms in the way sketched in [1.3], and refers for an elegant formulation to
\cite[Exercise~8, p.~48]{KreiselKrivine1967}; that referral is reproduced from
\cite[\S1.4, p.~899]{Statman1977}, and the cited exercise has not been
consulted here.  Kreisel and Tait proved \(K\) complete: if
\((E_1\land\cdots\land E_m)\to E_{m+1}\) is valid, then \(E_{m+1}\) is
derivable from \(E_1,\ldots,E_m\).  The chapter's aim is to answer (a), (b)
and (c) for \(K\), and in particular to show that direct computations are
formalizable in \(K\).

\begin{hbdefinition}{def:p02-w02-p02-028-st-J}{The axioms \(J\) [1.5]}
Consider the two-sorted first-order language of the structures
\(\langle\mathbb N,F,=,0,s\rangle\), where \(\mathbb N\) is the set of natural
numbers and \(F\) is a collection of number-theoretic functions that contains
the constant function \(0\) and the successor function \(s\) and is closed
under explicit definition (lambda-abstraction).  The axioms \(J\) are
\cite[\S1.5, p.~900]{Statman1977}:
\begin{enumerate}
  \item[(1)] \(x=x\),
  \item[(2)] \((fx=gx\land x=y)\to fy=gy\),
  \item[(3)] \((fx=gx\land y=x)\to fy=gy\),
  \item[(4)] \(sx=sy\to x=y\),
  \item[(5)] \(0=sx\to y=z\),
  \item[(\(6_n\))] \(x=s^nx\to y=z\),
\end{enumerate}
where \(s^1=s\) and \(s^{k+1}=ss^k\).
\end{hbdefinition}

The chapter records three facts about \(J\): its members are valid over all
structures \(\langle\mathbb N,F,=,0,s\rangle\); if
\((E_1\land\cdots\land E_m)\to E_{m+1}\) is not valid over the general models
of \(J\), then it has a countermodel of the form
\(\langle\mathbb N,F,=,0,s\rangle\); and the function variables of the
\(E_i\) can be interpreted by number-theoretic functions with finite support,
which gives decidability as well as completeness for \(K\)
\cite[\S1.5, p.~900]{Statman1977}.  This is the first, model-theoretic, proof
of Kreisel and Tait.  Theorem~\ref{thm:p02-w02-p02-028-st-J-route} proves
these facts under the hypotheses it states.

\emph{[1.6] The second proof and its by-products.}  The second proof
formalizes in \(K\) an elementary decision procedure for the validity of
\((E_1\land\cdots\land E_m)\to E_{m+1}\).  The result is a complete
proof-search procedure, and the analysis answers (a) and (c)
\cite[\S1.6, p.~900]{Statman1977}:
\begin{enumerate}
  \item[(a)] if \((E_1\land\cdots\land E_m)\to E_{m+1}\) is valid and
        \(E_1\land\cdots\land E_m\) is satisfiable, then there is a direct
        computation of \(E_{m+1}\) from \(E_1,\ldots,E_m\) of length
        \(\leq2^{kp}\), where \(p\) is the number of occurrences of
        non-logical symbols in the implication; if the conjunction is not
        satisfiable, a similar result holds, but not necessarily for
        \(E_{m+1}\);
  \item[(c)] validity of the implication can be decided by a deterministic
        Turing machine in time polynomial in the length of a code of the
        formula.
\end{enumerate}
The bound \(2^{kp}\) is for computations as trees of equation occurrences;
identifying occurrences of the same equation and coding computations as
sequences gives a bound of \(kp^2\).  The chapter presents the second proof
only.  The constant \(k\) is fixed by
Proposition~\ref{thm:p02-w02-p02-028-st-prop-3-2}; the claims are proved in
Theorems~\ref{thm:p02-w02-p02-028-st-completeness}
and~\ref{thm:p02-w02-p02-028-st-decision}.

\emph{[1.7] Question (b).}  For (b) the chapter analyses computations as in
Gentzen's cut-elimination argument: once individual inferences are fully
analysed, they can be permuted into a direct computation.  Precisely, it
announces that if in a computation \(T\) of \(E_{m+1}\) from
\(E_1,\ldots,E_m\)
\begin{enumerate}
  \item[(i)] the rules corresponding to axioms (5) and (\(6_n\)) are not
        used, and
  \item[(ii)] every application of substitution of equals for equals (the
        rule corresponding to axioms (2) and (3)) replaces at most one
        occurrence of one of the equals,
\end{enumerate}
then there is a direct computation of \(E_{m+1}\) from \(E_1,\ldots,E_m\) of
length \(\leq kp^2\), where \(p\) is the length of \(T\)
\cite[\S1.7, p.~901]{Statman1977}.  The assembly of this bound from
Propositions~\ref{thm:p02-w02-p02-028-st-prop-3-3-2},
\ref{thm:p02-w02-p02-028-st-prop-3-4} and~\ref{thm:p02-w02-p02-028-st-prop-3-5}
is Corollary~\ref{thm:p02-w02-p02-028-st-bound-b}.

\subsection{Preliminaries: terms, measures and computations}
\label{sec:p02-w02-p02-028-st-preliminaries}

\begin{hbdefinition}{def:p02-w02-p02-028-st-terms}{Equations, sets and substitution [2.1, 2.2]}
Equations \(E\) hold between individual terms \(a,b,c,\ldots\), which may
contain function variables; \(S\) ranges over finite sets of equations.
\([\cdots](a_1/x_1,\ldots,a_n/x_n)\) is the simultaneous substitution of
\(a_i\) for \(x_i\) in \([\cdots]\).  A substitution \(\theta\) may also
replace an \(n\)-place function variable \(f\) by a function term
\(\lambda x_1\cdots x_n a\), under the clauses
\[
 \theta(fa_1\cdots a_n)=(\theta f)(\theta a_1)\cdots(\theta a_n),
\]
\[
 (\lambda x_1\cdots x_na)b_1\cdots b_n=a(b_1/x_1,\ldots,b_n/x_n)
\]
\cite[\S\S2.1--2.2, p.~901]{Statman1977}.
\end{hbdefinition}

\begin{hbdefinition}{def:p02-w02-p02-028-st-measure}{Measure [2.3]}
A function \(M\) from terms to non-negative integers is a \emph{measure} if
\(M(a)\leq M(b)\) implies \(M(c(a/x))\leq M(c(b/x))\), and
\(M(a)\leq M(c(a/x))\) whenever \(x\) occurs in \(c\).  For a measure \(M\),
\(M(a=b)=M(a)+M(b)\) and \(M(S)=\sum_{E\in S}M(E)\).  Length is a measure,
where \(\operatorname{lh}(a)\) is the number of occurrences of symbols in
\(a\) \cite[\S2.3, p.~901]{Statman1977}.
\end{hbdefinition}

\begin{hbdefinition}{def:p02-w02-p02-028-st-nonprojecting}{Non-projecting substitution [2.4]}
A substitution \(\theta\) is \emph{non-projecting} if, whenever
\(\theta f=\lambda x_1\cdots x_na\), each \(x_i\) actually occurs in \(a\)
\cite[\S2.4, p.~901]{Statman1977}.
\end{hbdefinition}

\begin{hblemma}{thm:p02-w02-p02-028-st-measure-transfer}{proved=annex:p02-w02-p02-028-st-measure-transfer}{Measures under non-projecting substitution [2.4]}
If \(M\) is a measure and \(\theta\) is non-projecting, then
\(M^*(a)=M(\theta a)\) defines a measure.  The hypothesis cannot be dropped:
for the projecting \(\theta f=\lambda x\,0\), take \(c=fx\) and \(a=sx\); then
\(x\) occurs in \(c\), but \(\operatorname{lh}(\theta a)=2\) exceeds
\(\operatorname{lh}(\theta(c(a/x)))=\operatorname{lh}(0)=1\).
\end{hblemma}

\begin{hbdefinition}{def:p02-w02-p02-028-st-simple}{Simple set [2.5]}
\(S\) is \emph{simple} if each of its equations has one of the forms
\(x=0\), \(x_1=x_2\), \(x_1=sx_2\), or \(x_{m+1}=fx_1\cdots x_m\), with
individual variables \(x,x_i\) \cite[\S2.5, p.~901]{Statman1977}.
\end{hbdefinition}

\begin{hbdefinition}{def:p02-w02-p02-028-st-K}{The calculi \(K\) and \(H\) [2.6]}
The calculus \(K\) of Kreisel and Tait has the axioms \(a=a\) and four rules
\cite[\S2.6, p.~902]{Statman1977}: substitution of equals for equals
\[
 (1)\qquad \frac{E(a/x)\qquad a\doteq b}{E(b/x)},
\]
where \(a\doteq b\) stands ambiguously for \(a=b\) and \(b=a\); and
\[
 (2)\quad\frac{sa=sb}{a=b},
 \qquad
 (3)\quad\frac{0=sa}{b=c},
 \qquad
 (4_n)\quad\frac{a=s^na}{b=c}.
\]
The calculus with rule (1) only is \(H\); it corresponds to axioms (1), (2)
and (3) of \(J\).  Rule (2) corresponds to axiom (4), rule (3) to axiom (5),
and rule (\(4_n\)) to axiom (\(6_n\)).
\end{hbdefinition}

\begin{hbdefinition}{def:p02-w02-p02-028-st-computation}{Computations, length and singularity [2.7--2.9]}
A \emph{computation} \(T\) in \(K\) or \(H\) is a finite tree of equation
occurrences, each node with at most two premises, built from assumptions and
axioms by the rules; \(T\) is a computation of \(E\) from \(S\) when its root
is \(E\) and its assumptions belong to \(S\).  The \emph{length}
\(\operatorname{lh}(T)\) is the number of occurrences of equations in \(T\).
\(T\) is \emph{singular} if in each of its inferences of form (1) the
variable \(x\) occurs exactly once in \(E\)
\cite[\S\S2.7--2.9, p.~902]{Statman1977}.  The printed chapter numbers both
the length and the singularity definitions as 2.8; the second is the item
between 2.8 and 2.10.
\end{hbdefinition}

\begin{hbdefinition}{def:p02-w02-p02-028-st-moves}{Switches, shifts, moves and reduction [2.10]}
Given a computation \(T\), a \emph{switch} replaces a subcomputation of the
first form in Figure~\ref{fig:p02-w02-p02-028-st-moves} by the second, and a
\emph{shift} replaces a subcomputation of the third or the fifth form by the
fourth or the sixth, respectively.  A \emph{move} is a switch or a shift.
\(T_1\) \emph{reduces to} \(T_2\) if \(T_2\) results from \(T_1\) by a
finite sequence of moves \cite[\S2.10, pp.~902--903]{Statman1977}.
\end{hbdefinition}

\begin{figure}[htbp]
\centering
\small
\[
\begin{array}{c@{\qquad\text{by}\qquad}c}
\dfrac{\dfrac{E(a/x,b/y)\quad a\doteq c}{E(c/x,b/y)}\quad b\doteq d}{E(c/x,d/y)}
&
\dfrac{\dfrac{E(a/x,b/y)\quad b\doteq d}{E(a/x,d/y)}\quad a\doteq c}{E(c/x,d/y)}
\\[5ex]
\dfrac{\dfrac{E(a/x)\quad a\doteq c(b/y)}{E(c(b/y)/x)}\quad b\doteq d}{E(c(d/y)/x)}
&
\dfrac{E(a/x)\quad\dfrac{a\doteq c(b/y)\quad b\doteq d}{a\doteq c(d/y)}}{E(c(d/y)/x)}
\\[5ex]
\dfrac{\dfrac{E(b(a/x)/y)\quad a\doteq c}{E(b(c/x)/y)}\quad b(c/x)\doteq d}{E(d/y)}
&
\dfrac{E(b(a/x)/y)\quad\dfrac{b(c/x)\doteq d\quad c\doteq a}{b(a/x)\doteq d}}{E(d/y)}
\end{array}
\]
\caption{The moves of [2.10] (after \cite[pp.~902--903]{Statman1977}).  Row 1:
a switch.  Rows 2 and 3: the two shapes of shift.  In every shape the
premises are the conclusions of subcomputations \(T_1\) (left premise of the
upper inference), \(T_2\) (its right premise) and \(T_3\) (the right premise
of the lower inference); the replacement uses the same three subcomputations.}
\label{fig:p02-w02-p02-028-st-moves}
\end{figure}

\begin{hbdefinition}{def:p02-w02-p02-028-st-positions}{Positions, replacements, spines and \(H\) fragments}
A computation records, for each inference of form (1), the equation \(E\)
and the variable \(x\) of the rule; its left premise \(E(a/x)\) is the
\emph{main} premise and its right premise \(a\doteq b\) the \emph{side}
premise.  Positions in an equation are finite sequences of positive
integers: the left and right sides of \(l=r\) are at positions \(1\) and
\(2\), and the \(i\)-th argument of a term with outer symbol \(f\) or \(s\)
at position \(\pi\) is at position \(\pi i\).  Two positions are
\emph{comparable} if one is an initial segment of the other, and
\emph{disjoint} otherwise.  When \(x\) occurs in \(E\) exactly once, at
position \(\pi\), the inference has \(a\) at \(\pi\) in its main premise and
its conclusion is the main premise with \(b\) put at \(\pi\); we call it a
\emph{replacement of \(a\) by \(b\) at \(\pi\)}.  Thus every inference of
form (1) of a singular computation is a replacement at one position, and
conversely a replacement at \(\pi\) is the inference of form (1) whose
\(E\) is the main premise with a variable not occurring in it put at \(\pi\).

Two inferences \(I\), \(J\) of form (1) are \emph{nested to the left} if the
conclusion of \(I\) is the main premise of \(J\), and \emph{nested to the
right} if it is the side premise of \(J\).  A \emph{spine} is a maximal
sequence \(J_1,\ldots,J_m\) (\(m\geq1\)) of inferences of form (1) in which
\(J_{i+1}\) and \(J_i\) are nested to the left for each \(i<m\); \(J_1\) is
its bottom and \(J_m\) its top, and the main premise of \(J_m\), which is not
the conclusion of an inference of form (1), is the \emph{top equation}.
Since an equation occurrence is the main premise of at most one inference,
every inference of form (1) lies on exactly one spine.  An \emph{\(H\)
fragment} of a computation is a maximal set of inferences of form (1)
connected by the two nesting relations, together with their premises and
conclusions; its \emph{leaves} are its premises that are not conclusions of
its inferences.
\end{hbdefinition}

\begin{hbdefinition}{def:p02-w02-p02-028-st-leftness}{Leftness, \(T_R\) and \(T_L\) [2.11, 2.12]}
The \emph{leftness} of a computation \(T\) is the tuple
\((k_1,\ldots,k_n)\) in which \(k_i\) is the length of the \(i\)-th maximal
path of \(T\), counted from left to right; leftness is compared
lexicographically.  For singular \(T\), the chapter states
\cite[\S2.11, p.~903]{Statman1977} that
\begin{enumerate}
  \item[(a)] there is \(T_R\) such that \(T\) reduces to \(T_R\) and every
        sequence of moves beginning with \(T_R\) consists of switches; and
  \item[(b)] there is \(T_L\) such that \(T_L\) reduces to \(T\) and every
        sequence of moves ending in \(T_L\) consists of switches.
\end{enumerate}
\(T_R\) is obtained by repeating, as long as possible, a sequence of switches
followed by a shift, all involving inferences on the leftmost path of a
subcomputation; each such sequence decreases leftness.  \(T_L\) is obtained
by reversing the procedure, without the switches, as often as possible.  The
chapter then states that for singular \(H\) computations every
subcomputation of \(T_R\) and \(T_L\) has the forms of
Figure~\ref{fig:p02-w02-p02-028-st-fragments}, and that for singular \(K\)
computations the same holds of the \(H\) fragments of \(T_R\) and \(T_L\)
\cite[\S2.12, pp.~903--904]{Statman1977}.
Theorem~\ref{thm:p02-w02-p02-028-st-normal-forms} proves these statements;
in the terms of Definition~\ref{def:p02-w02-p02-028-st-positions}, the
inferences on the leftmost path of a subcomputation that are of form (1)
belong to one spine.
\end{hbdefinition}

\begin{figure}[htbp]
\centering
\small
\[
\begin{array}{c@{\qquad\qquad}c}
\begin{array}{c}
\dfrac{E(a_1/x_1,\ldots,a_n/x_n)\quad a_1\doteq b_1}{E(b_1/x_1,a_2/x_2,\ldots,a_n/x_n)}\\[1ex]
\ddots\\[1ex]
\dfrac{E(b_1/x_1,\ldots,b_{n-1}/x_{n-1},a_n/x_n)\quad a_n\doteq b_n}{E(b_1/x_1,\ldots,b_n/x_n)}
\end{array}
&
\begin{array}{c}
\dfrac{E_1\quad E_2}{E_3}\\[1ex]
\ddots\\[1ex]
\dfrac{E_n\quad E_{n+1}}{E_{n+2}}
\end{array}
\end{array}
\]
\caption{Subcomputation shapes of [2.12] (after
\cite[p.~903]{Statman1977}).  Left: a subcomputation of \(T_R\), with the
right premise \(a_i\doteq b_i\) the conclusion of \(T_i\); the main premise
descends on the left.  Right: the shape of \(T_L\), where each conclusion
\(E_3,E_5,\ldots\) is the main (left) premise of the next inference and the
side premises \(E_2,E_4,\ldots,E_{n+1}\) are not conclusions of inferences of
form (1).}
\label{fig:p02-w02-p02-028-st-fragments}
\end{figure}

\begin{hblemma}{thm:p02-w02-p02-028-st-moves-basic}{proved=annex:p02-w02-p02-028-st-moves-basic}{Moves preserve length, conclusion and singularity [2.11(i)]}
If \(T_1\) reduces to \(T_2\), then \(T_1\) and \(T_2\) have the same
conclusion, the same assumptions and the same length.  Assume that in every
inference of form (1) the substitution variable is chosen distinct from all
variables of the terms substituted.  If \(T_1\) is singular and \(T_1\)
reduces to \(T_2\) or \(T_2\) reduces to \(T_1\), then \(T_2\) is singular.
\end{hblemma}

The chapter further states that each combination of two inferences nested to
the left is covered by a clause of the move definition, and each combination
nested to the right arises from a clause of the shift definition
\cite[\S2.11(ii)--(iii), p.~903]{Statman1977}.

\begin{hblemma}{thm:p02-w02-p02-028-st-nesting}{proved=annex:p02-w02-p02-028-st-nesting}{Nested inferences [2.11(ii), (iii)]}
Let \(T\) be singular, and let \(I\) and \(J\) be inferences of form (1)
of \(T\), replacements at positions \(\rho\) and \(\pi\).
\begin{enumerate}
  \item[(ii)] If \(I\) and \(J\) are nested to the left, the subcomputation
        ending in \(J\) has the left shape of a switch, with \(I\) as its
        upper inference, when \(\rho\) and \(\pi\) are disjoint, and the left
        shape of one of the two shifts of
        Figure~\ref{fig:p02-w02-p02-028-st-moves} when they are comparable.
  \item[(iii)] If \(I\) and \(J\) are nested to the right, the subcomputation
        ending in \(J\) has the right shape of one of the two shifts, with
        \(I\) as the inference above its right premise.
\end{enumerate}
\end{hblemma}

\begin{hbtheorem}{thm:p02-w02-p02-028-st-normal-forms}{proved=annex:p02-w02-p02-028-st-normal-forms}{The forms \(T_R\) and \(T_L\) [2.11(a), (b), 2.12]}
Let \(T\) be a singular computation in \(K\) or \(H\).
\begin{enumerate}
  \item There is a computation \(T_R\) such that \(T\) reduces to \(T_R\),
        obtained by repeating a sequence of switches followed by a shift on
        one spine as long as possible, each repetition decreasing leftness.
        In \(T_R\) the positions of the inferences of each spine are pairwise
        disjoint, and every sequence of moves beginning with \(T_R\)
        consists of switches.
  \item There is a computation \(T_L\) such that \(T_L\) reduces to \(T\),
        obtained by repeating the inverse of a shift as long as possible,
        each repetition increasing leftness.  In \(T_L\) no side premise of
        an inference of form (1) is the conclusion of an inference of form
        (1), and every sequence of moves ending in \(T_L\) consists of
        switches.
  \item Every spine of \(T_R\) has the left form of
        Figure~\ref{fig:p02-w02-p02-028-st-fragments}, with top equation
        \(E(a_1/x_1,\ldots,a_n/x_n)\), each \(x_i\) occurring once in \(E\).
        Every \(H\) fragment of \(T_L\) is a single spine, of the right form
        of that figure; if \(T\) is an \(H\) computation with at least one
        inference, \(T_L\) is a single spine whose side premises are axioms
        or assumptions.
\end{enumerate}
\end{hbtheorem}

\begin{hbdefinition}{def:p02-w02-p02-028-st-m-direct}{\(M\)-direct computation [2.13]}
If \(M\) is a measure, a computation \(T\) of \(E\) from \(S\) is
\emph{\(M\)-direct} if for every term \(b\) occurring in \(T\) there is a
term \(c\) occurring in \(E\) or in \(S\) with \(M(b)\leq M(c)\)
\cite[\S2.13, p.~904]{Statman1977}.  ``Direct'' in [1.6] means
\(\operatorname{lh}\)-direct.
\end{hbdefinition}

\begin{hbconvention}{def:p02-w02-p02-028-st-convention}{Intended interpretation [2.14]}
Individual variables range over the natural numbers, function variables
range over number-theoretic functions, and \(0\) and \(s\) denote zero and
the successor function \cite[\S2.14, p.~904]{Statman1977}.  ``Valid'' and
``satisfiable'' for equations and their conjunctions refer to this
interpretation.
\end{hbconvention}

\begin{hbdefinition}{def:p02-w02-p02-028-st-analytic}{Analytic implication [2.15]}
\(\bigwedge S\to E\) is \emph{analytic} if it is valid and \(\bigwedge S\) is
satisfiable \cite[\S2.15, p.~904]{Statman1977}.
\end{hbdefinition}

When \(\bigwedge S\) is not satisfiable, one cannot expect
\(\operatorname{lh}\)-direct computations from \(S\) of every \(E\), only of
some instance of the premise of rule (3) or (\(4_n\)).  The chapter's example
is
\[
 S_n=\{x_1=sx_n,\ x_2=sx_1,\ \ldots,\ x_n=sx_{n-1}\},
 \qquad E=(0=s0)
\]
\cite[\S2.15, p.~904]{Statman1977}; Proposition~\ref{thm:p02-w02-p02-028-st-cycle}
works it out.  Accordingly, from here on \(K\) denotes the calculus without
rules (3) and (\(4_n\)); in this sense \(K\) is complete: if
\(\bigwedge S\to E\) is analytic, \(E\) is derivable from \(S\)
(Theorem~\ref{thm:p02-w02-p02-028-st-completeness}).

\subsection{Direct computations}
\label{sec:p02-w02-p02-028-st-direct}

\begin{hbtheorem}{thm:p02-w02-p02-028-st-prop-3-1}{proved=annex:p02-w02-p02-028-st-prop-3-1}{Direct computations from simple sets [3.1]}
Suppose \(S\) is simple and \(M\) is a measure.
\begin{enumerate}
  \item[(i)] If \(\bigwedge S\to x=y\) is analytic, there is an \(M\)-direct
        computation of \(x=y\) from \(S\) of length \(\leq2^{3k}\), where
        \(k\) is the number of variables occurring in \(S\).
  \item[(ii)] If \(\bigwedge S\) is not satisfiable, there is an
        \(M\)-direct computation of \(x=s^px\) or of \(0=sx\) from \(S\) of
        length \(\leq2^{4k}\), where \(p\leq\operatorname{card}(S)\), \(k\) is
        the number of variables occurring in \(S\), and \(x\) occurs in \(S\).
        (The annex proof gives \(p\leq\) the number of equations of the form
        \(x_1=sx_2\) in \(S\).)
\end{enumerate}
The computations use only rules (1) and (2) \cite[\S3.1, pp.~904--906]{Statman1977}.
\end{hbtheorem}

\begin{hbmethod}{def:p02-w02-p02-028-st-closure}{The identification procedure in the proof of [3.1]}
Input: a simple \(S\) and a measure \(M\).  Stage 1: \(S_1=S\cup\{x=x:x\text{
occurs in }S\}\).  Stage \(n+1\):
\begin{enumerate}
  \item[Case 1.] If \(S_n\) contains \(a=c\) and \(b=c\) with \(a\) not
        \(b\) and \(M(a)\leq M(b)\), put \(S_{n+1}=S_n(a/b)\).
  \item[Case 2.] Otherwise, if \(S_n\) contains \(c=sa\) and \(c=sb\) with
        \(a\) not \(b\) and \(M(a)\leq M(b)\), put \(S_{n+1}=S_n(a/b)\).
\end{enumerate}
Stop when neither case applies; the last stage is \(n_0\).  The proof keeps,
for each stage, a computation \(T^*_n\) of \(b=a\) from \(S\), a collection
\(C_n\) of computations from \(S\) of the members of \(S_n\), and the image
\(x_n=y_n\) of the target \(x=y\) \cite[\S3.1, pp.~904--905]{Statman1977}.
In both cases \(a\) and \(b\) are variables, so every stage deletes one
variable, and \(n_0\leq k\).
\end{hbmethod}

\begin{hbtheorem}{thm:p02-w02-p02-028-st-prop-3-2}{proved=annex:p02-w02-p02-028-st-prop-3-2}{Direct computations from arbitrary sets [3.2]}
If \(M\) is a measure, then
\begin{enumerate}
  \item[(i)] if \(\bigwedge S\to a=b\) is analytic, there is an \(M\)-direct
        computation of \(a=b\) from \(S\) of length
        \(\leq2^{9(\operatorname{lh}(S)+\operatorname{lh}(a=b))}\); and
  \item[(ii)] if \(\bigwedge S\) is not satisfiable, there is an
        \(M\)-direct computation of \(a=s^pa\) or of \(0=sa\) from \(S\) of
        length \(\leq2^{12\operatorname{lh}(S)}\), where
        \(p\leq\operatorname{lh}(S)\) and \(a\) occurs in \(S\)
\end{enumerate}
\cite[\S3.2, pp.~906--907]{Statman1977}.  The proof flattens \(S\) into a
simple set with one new variable \(x_t\) for each occurrence \(t\) of a term
and applies Theorem~\ref{thm:p02-w02-p02-028-st-prop-3-1} for the measure
\(M^*(c)=M(c(\ldots,t/x_t,\ldots))\).  The printed proof says ``apply
proposition'' without a number; the proposition applied is [3.1].
\end{hbtheorem}

\begin{hbtheorem}{thm:p02-w02-p02-028-st-completeness}{proved=annex:p02-w02-p02-028-st-completeness}{Completeness of \(K\), second proof}
Let \(S\) be finite and \(E\) an equation.
\begin{enumerate}
  \item If \(\bigwedge S\to E\) is analytic, then \(E\) is derivable from
        \(S\) by rules (1) and (2), by an \(\operatorname{lh}\)-direct
        computation of length \(\leq2^{9p}\), where \(p\) is the number of
        occurrences of non-logical symbols in \(\bigwedge S\to E\).
  \item If \(\bigwedge S\to E\) is valid, then \(E\) is derivable from \(S\)
        in the full calculus with rules (1)--(\(4_n\)).
  \item Every rule of the full calculus preserves truth under every
        interpretation of Convention~\ref{def:p02-w02-p02-028-st-convention};
        hence \(E\) is derivable from \(S\) in the full calculus iff
        \(\bigwedge S\to E\) is valid.
  \item Coding a computation as a sequence of equations, each an axiom, an
        assumption or a conclusion of a rule from earlier members, the
        computations of item~1 can be given as sequences of at most
        \(16p^2\) equations.
\end{enumerate}
\end{hbtheorem}

\begin{hbtheorem}{thm:p02-w02-p02-028-st-decision}{proved=annex:p02-w02-p02-028-st-decision}{Polynomial-time decision [1.6(c)]}
There is a deterministic algorithm that decides validity of
\(\bigwedge S\to a=b\) using a number of elementary operations polynomial in
\(p=\operatorname{lh}(S)+\operatorname{lh}(a=b)\); it runs the identification
procedure of Method~\ref{def:p02-w02-p02-028-st-closure} on the flattened set
and inspects the final stage.  Implemented on a deterministic Turing machine
on a code of the formula, it runs in polynomial time.
\end{hbtheorem}

\begin{hbtheorem}{thm:p02-w02-p02-028-st-J-route}{proved=annex:p02-w02-p02-028-st-J-route}{The model-theoretic route [1.5]}
\begin{enumerate}
  \item Every member of \(J\) is valid over every structure
        \(\langle\mathbb N,F,=,0,s\rangle\).
  \item If \(\bigwedge S\to E\) is not valid, it has a countermodel in which
        every function variable is interpreted by a number-theoretic function
        with finite support (value \(0\) at all but finitely many arguments)
        and every individual variable by a number below
        \((k+1)^2\), where \(k\) is the number of variables of the flattened
        set; hence validity is decidable by inspecting finitely many
        interpretations.
  \item Read a \emph{general model of \(J\)} as a two-sorted structure
        satisfying \(J\) whose function sort contains, for every term \(c\)
        and variable \(x\), an element \(\lambda x\,c\) applying as
        \((\lambda x\,c)a=c(a/x)\) (the closure under explicit definition
        assumed of \(F\)).  If \(\bigwedge S\to E\) fails in such a general
        model, it fails in some \(\langle\mathbb N,F,=,0,s\rangle\) under
        Convention~\ref{def:p02-w02-p02-028-st-convention}.
\end{enumerate}
\end{hbtheorem}

\subsection{Singular computations and permutation}
\label{sec:p02-w02-p02-028-st-singular}

\begin{hbdefinition}{def:p02-w02-p02-028-st-subterm}{Weak subterm and subterm properties [3.3.1]}
A computation \(T\) of \(E\) from \(S\) has the \emph{weak subterm property}
if each inference of form (2) in \(T\) has \(sa\) and \(sb\) occurring in
\(S\).  It has the \emph{subterm property} if, in addition, each axiom
\(a=a\) of \(T\) has \(a\) occurring in \(E\) or \(S\)
\cite[\S3.3.1, p.~907]{Statman1977}.
\end{hbdefinition}

\begin{hbtheorem}{thm:p02-w02-p02-028-st-prop-3-3-2}{proved=annex:p02-w02-p02-028-st-prop-3-3-2}{From the subterm property to \(M\)-directness [3.3.2]}
If \(M\) is a measure and \(T\) is a singular computation with the subterm
property, then \(T\) reduces to an \(M\)-direct computation \(T_M\)
\cite[\S3.3.2, p.~907]{Statman1977}.
\end{hbtheorem}

\begin{hbtheorem}{thm:p02-w02-p02-028-st-prop-3-4}{proved=annex:p02-w02-p02-028-st-prop-3-4}{From the weak to the full subterm property [3.4]}
If \(T\) is a singular computation of \(E\) from \(S\) with the weak subterm
property, then there is a singular computation \(T_S\) of \(E\) from \(S\)
with the subterm property and \(\operatorname{lh}(T_S)\leq\operatorname{lh}(T)\)
\cite[\S3.4, pp.~907--908]{Statman1977}.
\end{hbtheorem}

\begin{hbtheorem}{thm:p02-w02-p02-028-st-prop-3-5}{proved=annex:p02-w02-p02-028-st-prop-3-5}{Obtaining the weak subterm property [3.5]}
If \(T\) is a singular computation of \(E\) from \(S\), then there is a
singular computation \(T_W\) of \(E\) from \(S\) with the weak subterm
property and
\(\operatorname{lh}(T_W)\leq(\operatorname{lh}(T))^2+\operatorname{lh}(T)+1\)
\cite[\S3.5, pp.~908--910]{Statman1977}.
\end{hbtheorem}

The annex proof follows the chapter's route, one inference of form (2) at a
time from the top down, and its count gives the slightly stronger bound
\(\operatorname{lh}(T_W)\leq(\operatorname{lh}(T))^2\).  The three
transformations above compose: a singular computation is first given the
weak subterm property, then the subterm property, and is then permuted into
an \(M\)-direct one.

\begin{hbcorollary}{thm:p02-w02-p02-028-st-bound-b}{proved=annex:p02-w02-p02-028-st-bound-b}{The bound for question (b) [1.7]}
If \(T\) is a computation of \(E\) from \(S\) of length \(p\) satisfying
conditions (i) and (ii) of [1.7], then there is an
\(\operatorname{lh}\)-direct computation of \(E\) from \(S\) of length
\(\leq p^2\); in particular the chapter's bound \(kp^2\) holds with
\(k=1\).
\end{hbcorollary}

\subsection{Cases worked from the chapter}
\label{sec:p02-w02-p02-028-st-cases}

\begin{hbtheorem}{thm:p02-w02-p02-028-st-semigroup}{proved=annex:p02-w02-p02-028-st-semigroup}{Cancellation-semigroup rules [1.3(ii)]}
Let \(R_1,\ldots,R_m\) be equations between semigroup words over finitely
many generators, terms being built with one binary product.  For words
\(w_1,w_2\), the following are equivalent:
\begin{enumerate}
  \item \(w_1=w_2\) holds in every cancellation semigroup satisfying
        \(R_1,\ldots,R_m\) under every interpretation of the generators;
  \item \(w_1=w_2\) is derivable from \(R_1,\ldots,R_m\), the associativity
        axioms \(((ab)c)=(a(bc))\) and the identity axioms \(a=a\) by
        substitution of equals for equals and the two cancellation rules of
        [1.3(ii)].
\end{enumerate}
\end{hbtheorem}

\begin{hbtheorem}{thm:p02-w02-p02-028-st-horn}{proved=annex:p02-w02-p02-028-st-horn}{Horn derivations for the class \(L\) [1.3(ii)]}
Read the prime formulas of [1.3(i)] as equality-free atoms, with any
equality sign in them an ordinary binary relation symbol governed only by the
listed \(A_i\) and \(I_j\).  If \(G\in L\) and \(\exists x_1\cdots x_nG\) is
valid, then there are finitely many substitutions \(\theta_1,\ldots,\theta_r\)
of ground terms for \(x_1,\ldots,x_n\) such that \(B_{m+1}\) is derivable
from \(B_1,\ldots,B_m\) and the axioms \(\theta_lA_i\) by the rules
\(\theta_lP^j_1\cdots\theta_lP^j_{r(j)}/\theta_lC^j\).
\end{hbtheorem}

The proof uses Theorem~\ref{thm:p02-w02-p02-028-existential-validity} and
forward chaining.  Under logical equality the reading of
prime formulas as independent atoms fails
(Warning~\ref{def:p02-w02-p02-028-equality-limit}); the semigroup case is
covered separately by Theorem~\ref{thm:p02-w02-p02-028-st-semigroup}.

\emph{The route through cut elimination.}  The chapter obtains the same
statement from Gentzen's cut-elimination theorem
\cite[\S1.3(ii), p.~899]{Statman1977}.  Along that route the substitutions
\(\theta\) are read off a cut-free sequent derivation of
\(\Rightarrow\exists x_1\cdots\exists x_nG\): they are the instances of
\(G\) whose implication is introduced on the right.  The calculus is the
equality-free defined-negation calculus
\(\mathsf{LK}_{\mathrm d}\) (\texttt{sheet:lk-defneg}), for which
Theorem~\ref{thm:p07-classical-cut-elimination} proves cut elimination; the
passage from validity to a derivation uses the completeness of canonical
\(\mathsf{LK}\) (\texttt{sheet:lk})
(Theorem~\ref{thm:p07-canonical-lk-completeness}) and the translation of
Lemma~\ref{thm:p02-w02-p02-028-st-lk-translation}.

\begin{hbconvention}{def:p02-w02-p02-028-st-horn-cut-setting}{Setting for the cut-elimination route}
The signature is equality-free, and the prime formulas are the atomic
formulas \(R(t_1,\ldots,t_a)\); \(\bot\) is not prime, and an equality sign
in a prime formula of [1.3(i)] is read as an ordinary binary relation symbol.
The variables \(x_1,\ldots,x_n\) are pairwise distinct, each \(I_j\) has
\(r(j)\geq1\), and \(p+q+m\geq1\), so that the antecedent of \(G\) is a
formula.  That antecedent is written \(E\); it is the conjunction of the
occurrences of \(A_1,\ldots,A_p,I_1,\ldots,I_q,B_1,\ldots,B_m\) in some fixed
bracketing, and these occurrences are the \emph{leaves} of \(E\).
A \emph{substitution} \(\theta\) assigns a term \(\theta(x_i)\) to each
\(x_i\) and fixes every other variable; for a quantifier-free \(\chi\),
\(\theta\chi\) replaces every occurrence of each \(x_i\) by \(\theta(x_i)\)
simultaneously.  Since no \(B_k\) contains an \(x_i\), \(\theta B_k=B_k\).
Write \(\operatorname{par}(G)\) for the set of variables of \(G\) other than
\(x_1,\ldots,x_n\).  The \(l\)-th \emph{tail} of \(\theta I_j\), for
\(1\leq l\leq r(j)\), is
\(\theta\bigl(P^j_l\to(\cdots(P^j_{r(j)}\to C^j)\cdots)\bigr)\); the first
tail is \(\theta I_j\) itself.
\end{hbconvention}

\begin{hbdefinition}{def:p02-w02-p02-028-st-horn-system}{Left formulas, right formulas and the systems \(\mathsf H(\Theta;\Gamma)\)}
In the setting of
Convention~\ref{def:p02-w02-p02-028-st-horn-cut-setting}:
\begin{enumerate}
  \item A \emph{left formula} is a prime formula, a formula \(\theta E'\)
        with \(E'\) a subformula of \(E\) that is a conjunction, or a tail of
        some \(\theta I_j\).
  \item A \emph{right formula} is a prime formula or a \emph{\(G\)-form}
        \(\exists w_1\cdots\exists w_s\,\theta G\), where \(s\geq0\),
        \(\theta\) is a substitution and no \(w_i\) lies in
        \(\operatorname{par}(G)\); the \(w_i\) need not be distinct.  Its
        \emph{matrix substitution} is \(\theta\).  The formula
        \(\exists x_1\cdots\exists x_nG\) is a \(G\)-form with the identity
        substitution.
  \item Let \(\Theta\) be a finite set of substitutions and \(\Gamma\) a
        finite multiset of left formulas.  The system
        \(\mathsf H(\Theta;\Gamma)\) has as axioms \(B_1,\ldots,B_m\), every
        \(\theta A_i\) with \(\theta\in\Theta\), every prime member of
        \(\Gamma\), and, for each member \(\theta E'\) of \(\Gamma\), every
        \(\theta A_i\) with \(A_i\) a leaf of \(E'\) and every \(B_k\) that is
        a leaf of \(E'\).  Its rules
        are
        \(\theta P^j_1\cdots\theta P^j_{r(j)}/\theta C^j\) for
        \(\theta\in\Theta\) and every \(j\); the same rule for each member
        \(\theta E'\) of \(\Gamma\) and each \(j\) with \(I_j\) a leaf of
        \(E'\); and
        \(\theta P^j_l\cdots\theta P^j_{r(j)}/\theta C^j\) for each member of
        \(\Gamma\) that is the \(l\)-th tail of \(\theta I_j\).  A
        derivation in \(\mathsf H(\Theta;\Gamma)\) is a finite sequence of
        prime formulas, each an axiom or the conclusion of a rule whose
        premises all occur earlier in the sequence.
  \item The \emph{targets} of a finite multiset \(\Delta\) of right formulas
        are its prime members, together with \(B_{m+1}\) when some member of
        \(\Delta\) is a \(G\)-form.
\end{enumerate}
\(\mathsf H(\Theta;\varnothing)\) is the system of [1.3(ii)] restricted to
the substitutions in \(\Theta\): assumptions \(B_1,\ldots,B_m\), axioms
\(\theta A_i\), rules \(\theta P^j_1\cdots\theta P^j_{r(j)}/\theta C^j\).
\end{hbdefinition}

\begin{hblemma}{thm:p02-w02-p02-028-st-gform}{proved=annex:p02-w02-p02-028-st-gform}{Instances of \(G\)-forms}
If \(\exists w_1\psi\) is a \(G\)-form with matrix substitution \(\theta\),
then for every term \(t\) the formula \(\psi[t/w_1]\), formed with the
capture-avoiding substitution of
Definition~\ref{def:p02-fol-capture-avoiding-substitution}, is a \(G\)-form.
Its matrix substitution \(\theta'\) has
\(\theta'(x_i)=\eta(\theta(x_i))\), where \(\eta\) is a map from variables
to terms that fixes every member of \(\operatorname{par}(G)\) and
\(\eta(\theta(x_i))\) replaces each variable \(u\) of \(\theta(x_i)\) by
\(\eta(u)\).
\end{hblemma}

\begin{hblemma}{thm:p02-w02-p02-028-st-lk-translation}{proved=annex:p02-w02-p02-028-st-lk-translation}{From canonical LK to the defined-negation calculus}
If a finite equality-free sequent \(\Gamma\Rightarrow\Delta\) is derivable
on \texttt{sheet:lk}, then
\(d(\Gamma)\Rightarrow d(\Delta)\) is derivable on
\texttt{sheet:lk-defneg}, with \(d\) the
translation of Definition~\ref{def:p07-defined-negation-translation}.
Cut may occur in the derivation produced.
\end{hblemma}

This is the forward half of
Theorem~\ref{thm:p07-lk-presentation-equivalence}; the annex proves it rule
by rule.

\begin{hblemma}{thm:p02-w02-p02-028-st-extraction}{proved=annex:p02-w02-p02-028-st-extraction}{Horn derivations from cut-free derivations}
Assume Convention~\ref{def:p02-w02-p02-028-st-horn-cut-setting}.  Let
\(\mathcal D\) be a derivation on
\texttt{sheet:lk-defneg} containing neither Cut
nor Mix, with endsequent \(\Gamma\Rightarrow\Delta\), where \(\Gamma\) is a
multiset of left formulas and \(\Delta\) a multiset of right formulas.
Then there is a finite set \(\Theta\) of substitutions such that some target
of \(\Delta\) is derivable in \(\mathsf H(\Theta;\Gamma)\).  Each member of
\(\Theta\) is the matrix substitution of a quantifier-free \(G\)-form that is
the principal formula of an inference \(\to R\) in \(\mathcal D\).
\end{hblemma}

\begin{hbtheorem}{thm:p02-w02-p02-028-st-horn-cut}{proved=annex:p02-w02-p02-028-st-horn-cut}{The Horn statement of [1.3(ii)] from cut elimination}
Assume Convention~\ref{def:p02-w02-p02-028-st-horn-cut-setting} and let
\(G\in L\).
\begin{enumerate}
  \item If \(\Rightarrow\exists x_1\cdots\exists x_nG\) is derivable on
        \texttt{sheet:lk-defneg}, Cut and Mix
        allowed, then there is a finite set \(\Theta\) of substitutions for
        \(x_1,\ldots,x_n\) such that \(B_{m+1}\) is derivable from
        \(B_1,\ldots,B_m\) and the axioms \(\theta A_i\),
        \(\theta\in\Theta\), by the rules
        \(\theta P^j_1\cdots\theta P^j_{r(j)}/\theta C^j\),
        \(\theta\in\Theta\).  A fortiori \(B_{m+1}\) is derivable from
        \(B_1,\ldots,B_m\) and the axioms \(\theta A_i\) for each
        substitution \(\theta\) by the rules
        \(\theta P^j_1\cdots\theta P^j_{r(j)}/\theta C^j\) for each
        substitution \(\theta\), which is the conclusion of [1.3(ii)].
  \item If \(\exists x_1\cdots\exists x_nG\) is valid, that is, true in
        every structure under every assignment, then
        \(\Rightarrow\exists x_1\cdots\exists x_nG\) is derivable on
        \texttt{sheet:lk-defneg}, so the
        conclusion of item~1 holds.  No form of the Axiom of Choice is used.
\end{enumerate}
\end{hbtheorem}

\emph{The hypotheses the route uses.}  Cut elimination is applied on
\(\mathsf{LK}_{\mathrm d}\), which has no equality; on the equality sheet
\texttt{sheet:lk-eq} Cut is not admissible
(Theorem~\ref{thm:p07-eqlk-cut-failure}), so the argument does not run with
logical equality and the equality-free reading of prime formulas is a
hypothesis of the route, not only of
Proposition~\ref{thm:p02-w02-p02-028-st-horn}.  That \(B_k\) contains no
\(x_i\) gives \(\theta B_k=B_k\), so every instance of \(G\) has the same
assumptions and the same goal.  That \(\bot\) is not prime excludes the
initial sequent \(\bot,\Gamma\Rightarrow\Delta\) from a derivation of
\(\Rightarrow\exists\bar xG\).  Unlike
Proposition~\ref{thm:p02-w02-p02-028-st-horn}, the theorem does not require
\(\exists\bar xG\) to be a sentence, and the terms of the \(\theta\) may
contain variables.  The set \(\Theta\) bears on question (a) of [1.3(iii)]:
it consists of the substitutions of the instances of \(G\) that a cut-free
derivation actually introduces.

\begin{hbexample}{def:p02-w02-p02-028-st-horn-cut-example}{Two instances read off a cut-free derivation}
Take \(n=1\), \(p=0\), \(q=1\), \(m=1\), \(I_1=Q(x)\to Q(fx)\),
\(B_1=Q(c)\), \(B_2=Q(ffc)\), so \(G=(I_1\land B_1)\to B_2\), and write
\(G_t=((Q(t)\to Q(ft))\land Q(c))\to Q(ffc)\) for the instance
\(x\mapsto t\).  A cut-free derivation of \(\Rightarrow\exists x\,G\):
\(\to L\) on \(Q(c)\to Q(fc)\) with the identities \(Q(c)\Rightarrow Q(c)\)
and \(Q(fc)\Rightarrow Q(fc)\) gives
\(Q(c)\to Q(fc),Q(c)\Rightarrow Q(fc)\); \(\to L\) on
\(Q(fc)\to Q(ffc)\) with \(Q(ffc)\Rightarrow Q(ffc)\) gives
\(Q(fc)\to Q(ffc),\ Q(c)\to Q(fc),\ Q(c)\Rightarrow Q(ffc)\).  Left
weakening by a second \(Q(c)\), right weakening by a second \(Q(ffc)\) and
two \(\land L\) give
\((Q(c)\to Q(fc))\land Q(c),\ (Q(fc)\to Q(ffc))\land Q(c)\Rightarrow
Q(ffc),Q(ffc)\).  Two \(\to R\) give \(\Rightarrow G_c,G_{fc}\); two
\(\exists R\) give \(\Rightarrow\exists xG,\exists xG\); right contraction
gives \(\Rightarrow\exists xG\).  The principal formulas of the two
\(\to R\) are \(G_{fc}\) and \(G_c\), so
Lemma~\ref{thm:p02-w02-p02-028-st-extraction} yields
\(\Theta=\{x\mapsto c,\ x\mapsto fc\}\), and the derivation in
\(\mathsf H(\Theta;\varnothing)\) is \(Q(c)\), \(Q(fc)\), \(Q(ffc)\): the
assumption \(B_1\), then the rule \(Q(c)/Q(fc)\), then the rule
\(Q(fc)/Q(ffc)\).

One substitution never suffices here.  For a term \(t\), take the structure
whose domain is the set of terms, with every term denoting itself under the
identity assignment, \(Q\) true of \(c\), also true of \(fc\) when \(t\) is
\(c\), and false elsewhere.  It makes \(Q(c)\) and \(Q(t)\to Q(ft)\) true
and \(Q(ffc)\) false, so \(G_t\) is not valid.  The right contraction of two
instances is therefore essential to the cut-free derivation.
\end{hbexample}

\begin{hbwarning}{def:p02-w02-p02-028-st-horn-cut-limits}{What the cut-elimination route does not give}
The route gives no bound on the number of substitutions in \(\Theta\) or on
the size of their terms in terms of \(G\): the cut-elimination construction
of Theorem~\ref{thm:p07-classical-cut-elimination} proves no size bound, and
no recursive bound on the number of instances exists
(Corollary~\ref{thm:p02-w02-p02-028-no-recursive-bound}).  It does not cover
logical equality, as explained above; the cancellation-semigroup statement of
[1.3(ii)] is the equality case, and
Theorem~\ref{thm:p02-w02-p02-028-st-semigroup} proves it from the
cancellation rules and substitution of equals for equals instead.
\end{hbwarning}

\begin{hbtheorem}{thm:p02-w02-p02-028-st-cycle}{proved=annex:p02-w02-p02-028-st-cycle}{The example \(S_n\) of [2.15]}
Let \(n\geq2\) and \(S_n=\{x_1=sx_n,\ x_2=sx_1,\ \ldots,\ x_n=sx_{n-1}\}\).
\begin{enumerate}
  \item \(\bigwedge S_n\) is not satisfiable, so \(\bigwedge S_n\to0=s0\) is
        valid but not analytic.
  \item \(x_1=s^nx_1\) has an \(\operatorname{lh}\)-direct computation from
        \(S_n\) by rule (1) for the measure \(\operatorname{lh}\) relative to
        its own conclusion, and \(0=s0\) follows by rule (\(4_n\)).
  \item No computation of \(0=s0\) from \(S_n\) in the full calculus with rules
        (1)--(\(4_m\)) is \(\operatorname{lh}\)-direct.
\end{enumerate}
For \(n=1\) the rule (\(4_1\)) applied to the assumption \(x_1=sx_1\) gives
an \(\operatorname{lh}\)-direct computation of \(0=s0\) of length \(2\).
\end{hbtheorem}

\begin{hbtheorem}{thm:p02-w02-p02-028-st-prop-4-1-1}{proved=annex:p02-w02-p02-028-st-prop-4-1-1}{Substitution instances and directness [4.1.1]}
If \(M\) is a measure and \(E\) is derivable from \(S\) in \(K\) (rules
(1) and (2), as fixed after [2.15]), then there are
substitution instances \(E_1,S_1;\ldots;E_n,S_n\) of \(E\) and \(S\) and
singular computations \(T_i\) of \(E_i\) from \(S_i\) such that for every
substitution \(\theta\) there are a non-projecting \(\theta_i\) and a
computation \(T_\theta\) with
(i) \(\theta E=\theta_iE_i\), (ii) \(\theta S=\theta_iS_i\),
(iii) \(T_i\) reduces to \(T_\theta\), and (iv) \(\theta_iT_\theta\) is
\(M\)-direct \cite[\S4.1.1, pp.~910--911]{Statman1977}.
\end{hbtheorem}

Here a substitution \(\theta\) replaces individual variables by terms and
function variables by function terms as in [2.2], and \(\theta T\) applies
\(\theta\) to every equation of \(T\).  The chapter draws from [4.1.1] the
consequence that the rules of \(K\) are, up to a constant factor for
sequence-coded computations, the most efficient of any finite set of
``analytic'' rules, in the length of computations and of terms, and leaves
most of the details to the reader \cite[\S4.1, p.~910]{Statman1977}.  The
chapter does not define an analytic rule.  The theorem below reads one as a
rule whose premises and conclusion form an analytic implication in the sense
of [2.15], and proves the claim under that reading.

\begin{hbtheorem}{thm:p02-w02-p02-028-st-optimality}{proved=annex:p02-w02-p02-028-st-optimality}{Efficiency of the rules of \(K\) [4.1]}
Let \(\mathcal R\) be a finite set of rules, each given by a finite set
\(S_\rho\) of equations and an equation \(E_\rho\) such that
\(\bigwedge S_\rho\to E_\rho\) is analytic; an application of the rule is an
inference from the members of \(\theta S_\rho\) to \(\theta E_\rho\), for a
substitution \(\theta\).  There is a number \(c\), depending only on
\(\mathcal R\), such that the following holds.  Let \(D\) be a sequence of
\(\ell\) equations ending in \(E\), each an axiom \(a=a\), a member of \(S\),
or the conclusion of an application of a rule of \(\mathcal R\) whose
premises occur earlier in \(D\).  Then there is a sequence of at most
\(c\ell\) equations ending in \(E\), each an axiom, a member of \(S\), or
the conclusion of an inference of form (1) or (2) whose premises occur
earlier, in which every term has length at most the length of some term
occurring in \(D\).
\end{hbtheorem}

\begin{hbwarning}{def:p02-w02-p02-028-st-open}{What the chapter leaves open [4.2, 4.3]}
The chapter states that its analysis does not settle the relation of arbitrary computations to direct ones, and poses the
\emph{Question}: is there an infinite sequence of pairs \(E_i,S_i\) and a
number \(n\) such that \(E_i\) has a computation from \(S_i\) of length
\(\leq n\) while every direct computation of \(E_i\) from \(S_i\) has length
at least \(i\)?  A negative answer would follow from the \emph{Statement}:
there is a function \(f\) such that for every computation \(T\) of \(E\)
from \(S\) there are \(E^*,S^*,\theta^*,T^*\) with (i) \(\theta^*\)
non-projecting, (ii) \(\theta^*E^*=E\), (iii) \(\theta^*S^*=S\), (iv)
\(T^*\) a singular computation of \(E^*\) from \(S^*\), and (v)
\(\operatorname{lh}(T^*)\leq f(\operatorname{lh}(T))\)
\cite[\S4.2, p.~911]{Statman1977}.  Neither the Question nor the Statement is
settled there.  Finally [4.3] points to \cite[Chapter~2]{Statman1974}, where
computations from axioms satisfiable by partial functions are compared in
length with computations from ``equivalent'' axioms satisfiable only by
total functions.
\end{hbwarning}

\emph{Relationships.}  [1.1] is the existential case of
Theorem~\ref{thm:p02-w02-p02-028-existential-validity}.  The identification
procedure of [3.1] merges variables that are forced equal, and the read-off
interpretation in its proof plays for equations the part that the Herbrand
structure of Lemma~\ref{thm:p02-w02-p02-028-herbrand-model} plays for atoms.
Switches and shifts permute substitution inferences in the way Gentzen's
reductions permute rules, which is the sense of ``direct'' in the chapter's
title.

\emph{Sources.}  The chapter's reference list \cite[p.~912]{Statman1977}
contains \cite{KreiselKrivine1967}, \cite{KreiselTait1961},
\cite{Lifsic1968} and \cite{Statman1974}; the item by Lif\v{s}ic is listed
there without a citation in the text of the chapter.  The general Herbrand
results of Sections~\ref{sec:p02-w02-p02-028-herbrand-structures}
and~\ref{sec:p02-w02-p02-028-validity-forms} are proved from the handbook's own
definitions; for the existential case see also
\cite[Theorem~2.9, p.~876]{Schwichtenberg1977}.
